\chapter{数项级数}
\label{chap:numerical-series-solutions}

\begin{solution}{series-cauchy-one-term}
在 Cauchy 准则中取 \(m=n+1\)，得到
\(\abs{a_{n+1}}<\varepsilon\) 对充分大的 \(n\) 成立，故
\(a_n\to0\)。调和级数通项 \(1/n\to0\)，但按正文分组证明其部分和
无上界，所以逆命题失败。
\end{solution}

\begin{solution}{series-geometric-complete}
若 \(q\ne1\)，
\[
 s_n=\sum_{k=0}^nq^k=\frac{1-q^{n+1}}{1-q}.
\]
当 \(\abs q<1\) 时 \(q^{n+1}\to0\)，总和为 \(1/(1-q)\)。
\(q=1\) 时部分和为 \(n+1\)；\(q=-1\) 时部分和振荡；\(\abs q>1\)
时通项不趋零。故恰在 \(\abs q<1\) 时收敛。
\end{solution}

\begin{solution}{series-positive-bounded-partial}
正项级数的部分和单调不减。它收敛则当然有界；若有上界，单调有界收敛
定理给出有限极限，故级数收敛。
\end{solution}

\begin{solution}{series-finite-change}
若只改动前 \(N\) 项，则从第 \(N\) 个部分和以后，新旧部分和相差常数
\[
 C=\sum_{k=1}^N(\widetilde a_k-a_k).
\]
一个部分和数列收敛当且仅当另一个收敛；收敛时新和等于旧和加 \(C\)。
\end{solution}

\begin{solution}{series-limit-comparison-zero-infinity}
若 \(a_n/b_n\to0\)，则尾部 \(a_n\le b_n\)，所以
\(\sum b_n\) 收敛推出 \(\sum a_n\) 收敛；不能反推，例如
\(a_n=1/n^2,b_n=1/n\)。若比值趋于无穷，则尾部 \(a_n\ge b_n\)，
所以 \(\sum b_n\) 发散推出 \(\sum a_n\) 发散；不能反推，例如
\(a_n=1/n,b_n=1/n^2\)。
\end{solution}

\begin{solution}{series-harmonic-blocks}
从 \(2^{k-1}+1\) 到 \(2^k\) 有 \(2^{k-1}\) 项，每项至少 \(2^{-k}\)，
故该块和至少 \(1/2\)。因此
\[
 s_{2^m}\ge1+\frac m2,
\]
随 \(m\to\infty\) 无界。
\end{solution}

\begin{solution}{series-root-factorial}
设 \(a_n=n!/n^n\)。则
\[
 \frac{a_{n+1}}{a_n}
 =\frac{(n+1)!}{(n+1)^{n+1}}\frac{n^n}{n!}
 =\left(\frac{n}{n+1}\right)^n\longrightarrow e^{-1}<1.
\]
比值判别给出收敛。
\end{solution}

\begin{solution}{series-root-polynomial-exponential}
\[
 \sqrt[n]{\abs{n^pr^n}}
 =\abs r\,n^{p/n}\longrightarrow\abs r.
\]
当 \(\abs r<1\) 时绝对收敛；当 \(\abs r>1\) 时通项不趋零。
若 \(\abs r=1\)，绝对值级数为 \(\sum n^p\)，由 \(p\)-级数判别，
它恰在 \(p<-1\) 时收敛。\(r=0\) 时除有限项外全为零。
\end{solution}

\begin{solution}{series-ratio-inconclusive-p}
对 \(a_n=n^{-p}\)，
\[
 \frac{a_{n+1}}{a_n}=\left(\frac n{n+1}\right)^p\to1,
\qquad
 \sqrt[n]{a_n}=n^{-p/n}\to1.
\]
\(p=1\) 时级数发散，\(p=2\) 时收敛，所以临界值 \(1\) 不含足够信息。
\end{solution}

\begin{solution}{series-integral-remainder-estimate}
对单调递减的 \(f(x)=x^{-p}\)，逐区间比较给出
\[
 \int_{N+1}^\infty x^{-p}dx
 \le\sum_{n=N+1}^\infty n^{-p}
 \le\int_N^\infty x^{-p}dx.
\]
计算两端积分即得题式。
\end{solution}

\begin{solution}{series-condensation-log}
对
\[
 a_n=\frac1{n\log n(\log\log n)^q}
\]
从充分大指标起各因子为正且 \(a_n\) 单调递减，可以使用凝聚判别。凝聚项为
\[
 2^k a_{2^k}
 =\frac{1}{k\log2\,\bigl(\log(k\log2)\bigr)^q}.
\]
由于 \(\log(k\log2)/\log k\to1\)，极限比较把它化为常数倍的
\[
 \frac{1}{k(\log k)^q}
\]
。再次使用积分判别：令 \(t=\log x\)，则
\[
 \int^\infty\frac{dx}{x(\log x)^q}
 =\int^\infty t^{-q}\,dt,
\]
恰在 \(q>1\) 时收敛。因此原级数也恰在 \(q>1\) 时收敛。
\end{solution}

\begin{solution}{series-raabe-log-critical}
计算
\[
 \frac{a_n}{a_{n+1}}
 =\frac{n+1}{n}
  \left(\frac{\log(n+1)}{\log n}\right)^q
 =1+\frac1n+\frac{q}{n\log n}
  +o\!\left(\frac1{n\log n}\right).
\]
所以 Raabe 表达式趋于 \(1\)，与 \(q\) 无关；但该级数在
\(q>1\) 时收敛、\(q\le1\) 时发散，说明临界值必须保留更低阶项。
\end{solution}

\begin{solution}{series-gauss-gamma-ratio}
题设给出
\[
 \frac{a_n}{a_{n+1}}=1+\frac{\alpha}{n}+O(n^{-2}).
\]
取对数时 \(\log(1+y)=y+O(y^2)\)，故
\[
 \log a_{n+1}-\log a_n
 =-\frac{\alpha}{n}+O(n^{-2}).
\]
从固定指标求和到 \(n-1\)，利用
正文引理 \ref{lem:harmonic-log-correction} 的
\(\sum_{k<n}1/k=\log n+\gamma+o(1)\) 和
\(\sum k^{-2}<\infty\)，可知
\[
 \log a_n+\alpha\log n\longrightarrow\ell
\]
对某个实数 \(\ell\) 成立。于是 \(a_nn^\alpha\to e^\ell=:C>0\)，
即 \(a_n\sim Cn^{-\alpha}\)。与 \(p\)-级数作极限比较可得：
\(\alpha>1\) 时收敛，\(\alpha\le1\) 时发散；特别地，
\(\alpha=1\) 时与调和级数等价。
\end{solution}

\begin{solution}{series-sparse-harmonic}
非零项只在 \(n=2^k\) 出现，故原级数等于
\[
 \sum_{k=0}^\infty\frac1{2^k}=2.
\]
分母集合 \(\{2^k:k\ge0\}\) 显然无限。
\end{solution}

\begin{solution}{series-divergent-slow-terms}
令
\[
 a_n=\begin{cases}
 n^{-2},&n=2^k,\\
 n^{-1},&\text{其他 }n.
 \end{cases}
\]
取 \(n_k=2^k\)，则 \(n_ka_{n_k}=2^{-k}\to0\)，而整个数列也趋于零。
从调和级数中删去的项之和
\(\sum_k2^{-k}\) 收敛，所以剩余调和项仍发散；加上正项
\(\sum 2^{-2k}\) 后仍发散。
\end{solution}

\begin{solution}{series-comparison-oscillation-positive}
取 \(b_n=1/n^2\)，并令
\[
 a_n=\begin{cases}
 b_n,&n\text{ 偶},\\
 2b_n,&n\text{ 奇}.
 \end{cases}
\]
则 \(a_n/b_n\) 在 \(1,2\) 之间振荡而无极限，但
\[
 b_n\le a_n\le2b_n,
\]
所以两级数都收敛。把 \(b_n\) 改为 \(1/n\) 可得同样结构的发散例。
\end{solution}

\begin{solution}{series-cauchy-condensation-base}
对块 \(m^k\le n<m^{k+1}\)，其中共有
\((m-1)m^k\) 个整数。单调性给出
\[
 (m^{k+1}-m^k)a_{m^{k+1}}
 \le\sum_{n=m^k}^{m^{k+1}-1}a_n
 \le(m^{k+1}-m^k)a_{m^k}.
\]
右端表明 \(\sum m^ka_{m^k}\) 收敛时原级数收敛。左端改写为
\[
 (m-1)m^k a_{m^{k+1}}
 =\frac{m-1}{m}\,m^{k+1}a_{m^{k+1}},
\]
只差一个正常数和一次指标平移，所以原级数收敛时凝聚级数也收敛。
\end{solution}

\begin{solution}{series-kummer-test}
第一种条件乘以 \(a_{n+1}\) 得
\[
 p_na_n-p_{n+1}a_{n+1}\ge c\,a_{n+1}.
\]
从 \(N\) 到 \(m\) 求和，左端望远镜消去，故
\[
 c\sum_{n=N+1}^{m+1}a_n\le p_Na_N.
\]
正项部分和有界，级数收敛。

若表达式不大于零，则
\[
 p_na_n\le p_{n+1}a_{n+1}.
\]
所以 \(p_na_n\ge p_Na_N>0\)，从而
\(a_n\ge C/p_n\)。若 \(\sum1/p_n\) 发散，比较判别给出
\(\sum a_n\) 发散。

取 \(p_n=n\)，Kummer 表达式为
\[
 n\left(\frac{a_n}{a_{n+1}}-1\right)-1.
\]
其最终正下界与最终非正情形分别给出 Raabe 判别的两侧。
\end{solution}
